\section{Анализ предметной области и постановка задачи}

\subsection{Уведомление как распределённый процесс}

В работе уведомление трактуется как управляемый распределённый процесс, а не как одиночный вызов внешнего email- или SMS-провайдера. Для прикладной системы существенен не сам факт отправки сообщения во внешний канал, а полный маршрут обработки: от приёма доменной команды до появления итогового статуса, по которому можно восстановить результат доставки.

Начальной точкой этого маршрута служит уведомительное событие. Под уведомительным событием понимается зарегистрированное намерение платформы выполнить доставку сообщения определённой аудитории с заданным шаблоном, набором параметров, приоритетом и каналом. Событие фиксирует бизнес-смысл уведомления и отделяет его от конкретной попытки фактической отправки.

Следующий уровень модели образует запуск доставки. Один и тот же экземпляр уведомительного события может быть запущен сразу после создания либо в отложенном режиме. Запуск фиксирует момент исполнения, состав аудитории и параметры обработки, то есть переводит намерение в исполняемое действие.

Наиболее прикладной стадией является канальная доставка. Она выполняется для конкретного получателя и конкретного канала связи, например email или SMS. Именно здесь принимается решение о допустимости отправки, выполняется обращение к внешнему провайдеру, фиксируются временные ошибки, повторные попытки, окончательный результат либо осознанный пропуск.

Завершающим слоем модели выступает история доставки. Она хранит агрегированную последовательность статусов и позволяет связать входную команду, запуск, канальную обработку и итоговое состояние в один наблюдаемый маршрут. Поэтому далее уведомление рассматривается как цепочка ``уведомительное событие --- запуск доставки --- канальная доставка --- история'', а не как изолированная операция отправки.

\subsection{Гарантированность доставки и границы ответственности}

Под гарантированной доставкой в рамках данной работы понимается не физическая гарантия прочтения сообщения конечным пользователем и не однократная безусловная доставка внешнему провайдеру. Внешние канальные провайдеры рассматриваются как ненадёжные компоненты, способные возвращать ошибки, таймауты и неоднозначные результаты.

Внутри платформы обеспечиваются следующие свойства обработки:
\begin{enumerate}
    \item если команда на создание уведомительного события успешно принята фасадом и зафиксирована в Facade PostgreSQL, событие сохраняется и может быть опубликовано после сбоя отдельного экземпляра сервиса;
    \item переход от сохранённого события к асинхронной обработке выполняется через transactional outbox, когда доменное изменение и запись интеграционного сообщения об уведомлении фиксируются в одной локальной транзакции, позволяя гарантировать атомарность записи на уровне базы данных;
    \item повторная публикация или повторное получение сообщения допускается, но не должна приводить к неконтролируемому созданию дублей благодаря idempotency key, transactional inbox и уникальным ограничениям;
    \item временные ошибки обработки переводятся в контролируемые повторные попытки, а окончательные ошибки фиксируются в статусной модели;
    \item итог обработки сохраняется в истории доставки и может быть восстановлен по техническим и доменным идентификаторам.
\end{enumerate}

Гарантия относится к сохранности, повторяемости и трассируемости обработки внутри платформы. Фактическая доставка конечному пользователю зависит от внешнего провайдера и подтверждается только статусами, полученными или сформированными канальным сервисом.

\subsection{Мультиканальность платформы уведомлений}

Под мультиканальностью в рамках работы понимается способность платформы обрабатывать уведомительные события через несколько независимых канальных контуров при сохранении единой доменной модели, единого механизма запуска доставки и единого контура истории.

В текущей реализации поддержаны два канала: email и SMS. Каждый канал представлен отдельным sender-сервисом, имеет собственный контур обработки, собственные попытки отправки и собственные интеграционные особенности. При этом оба канала используют общие архитектурные механизмы: уведомительное событие, запуск доставки, Kafka-публикацию, проверку профиля получателя, retry, фиксацию статуса и передачу результата в history-writer.

Добавление нового канала предполагает создание нового канального обработчика и расширение набора канальных представлений шаблона, но не требует изменения базовой модели уведомительного события, запуска доставки и истории.

\subsection{Проблемы прямого синхронного подхода}

Подход, при котором клиентский запрос удерживается до момента фактической отправки сообщения внешнему провайдеру, создаёт для платформы уведомлений несколько прикладных ограничений:
\begin{enumerate}
    \item сбой внешних канальных провайдеров напрямую ухудшает время ответа и повышает вероятность таймаутов клиентского запроса;
    \item повторные попытки в условиях сетевой неопределённости создают риск дублей без явной идемпотентности;
    \item приём команд, планирование, проверка профиля и канальная доставка оказываются жёстко связаны и не масштабируются независимо;
    \item временные и окончательные ошибки трудно различать в линейном сценарии;
    \item канальные сервисы начинают дублировать логику проверки профиля получателя и правил маршрутизации.
\end{enumerate}

Для предметной области уведомлений это означает, что прямой синхронный подход удобен только на малом масштабе и при отсутствии требований к управляемой повторной обработке.

\subsection{Анализ аналогов и типовых подходов}

До выбора внутренней архитектуры были рассмотрены типовые способы организации уведомлений, применяемые в прикладных системах. Сравнение проводилось не по общему удобству разработки, а по тому, насколько подход позволяет управлять полным жизненным циклом уведомительного события внутри системы.

\begin{table}[H]
    \centering
    \footnotesize
    \renewcommand{\arraystretch}{1.15}
    \caption{Сравнение типовых подходов к организации доставки уведомлений}
    \label{tab:notification-approaches}
    \begin{tabular}{|L{0.22\textwidth}|L{0.20\textwidth}|L{0.27\textwidth}|L{0.21\textwidth}|}
        \hline
        \textbf{Подход} & \textbf{Сильные стороны} & \textbf{Ограничения} & \textbf{Вывод для работы} \\
        \hline
        Прямая интеграция каждого сервиса с email/SMS-провайдером & Минимальная начальная сложность, быстрый запуск одного сценария & Дублирование retry и idempotency, отсутствие единой истории, расхождение правил профиля получателя, сложность расследования инцидентов & Не выбран как целевой подход \\
        \hline
        Внешняя SaaS notification-платформа & Готовые интеграции с несколькими каналами, меньше собственной инфраструктуры & Зависимость от внешнего поставщика, ограничения на внутренние данные и жизненный цикл события & Может использоваться как внешний provider, но не как внутренний контур управления \\
        \hline
        Очередь задач и worker & Проще, чем событийная платформа, выводит отправку из синхронного запроса & Слабее история и повторное чтение, хуже разделяются событие, запуск и канальная доставка & Подходит для простых сценариев \\
        \hline
        Внутренняя событийная платформа с Kafka и outbox/inbox & Отделяет командный контур от доставки, поддерживает независимые sender-сервисы, историю и безопасную повторную обработку & Выше начальная сложность, требуется дисциплина контрактов, метрик и дедупликации & Выбранный подход \\
        \hline
    \end{tabular}
\end{table}

Прямая интеграция каждого сервиса с внешним email- или SMS-провайдером не удовлетворяет требованиям централизованной истории, единых правил профиля получателя и управляемой повторной обработки. Очередь задач и worker выводят отправку из синхронного запроса, но хуже подходят для маршрута, в котором требуется раздельно фиксировать событие, запуск доставки, канальную обработку и статусную историю. Внешняя SaaS-платформа может быть полезна как канал или шлюз, однако она не решает задачу внутреннего управления уведомительным событием и не заменяет собственный контур истории и профиля.

Для данной ВКР выбран внутренний событийный подход, поскольку предметом работы является не только отправка сообщения во внешний канал, а управление полным жизненным циклом уведомительного события: приёмом команды, запуском доставки, канальной обработкой и фиксацией результата.

\subsection{Требования к платформе}

Функциональные требования сводятся к следующим возможностям:
\begin{enumerate}
    \item создание, изменение, отмена, чтение и постраничный просмотр уведомительных событий;
    \item управление аудиторией события;
    \item запуск доставки в немедленном или отложенном режиме;
    \item подготовка содержимого уведомления на основе шаблона и параметров полезной нагрузки;
    \item централизованная проверка профиля получателя и допустимости отправки;
    \item передача запуска доставки в асинхронный контур обработки;
    \item фиксация статусов и истории попыток доставки;
    \item чтение агрегированной истории по получателю и событию;
    \item расширение платформы на новые каналы без переработки базовой доменной модели.
\end{enumerate}

Нефункциональные требования приведены в таблице~\ref{tab:nfr-platform}.

\begin{table}[H]
    \centering
    \footnotesize
    \renewcommand{\arraystretch}{1.15}
    \caption{Нефункциональные требования платформы уведомлений}
    \label{tab:nfr-platform}
    \begin{tabular}{|L{0.24\textwidth}|L{0.68\textwidth}|}
        \hline
        \textbf{Категория} & \textbf{Требование} \\
        \hline
        Сохранность события & Принятая команда не должна теряться при частичном отказе отдельного сервиса или экземпляра. \\
        \hline
        Идемпотентность & Повторный запрос или повторное получение сообщения не должны создавать неконтролируемые дубли. \\
        \hline
        Масштабируемость & Приём команд, планирование, канальная обработка и история должны масштабироваться независимо. \\
        \hline
        Наблюдаемость & По метрикам и журналам должен восстанавливаться маршрут от входной команды до записи статуса. \\
        \hline
        Расширяемость & Добавление нового канала не должно требовать изменения базовой модели уведомительного события и запуска доставки. \\
        \hline
    \end{tabular}
\end{table}

\begin{table}[H]
    \centering
    \footnotesize
    \renewcommand{\arraystretch}{1.15}
    \caption{Проверка выполнения ключевых требований}
    \label{tab:nfr-validation}
    \begin{tabular}{|L{0.34\textwidth}|L{0.54\textwidth}|}
        \hline
        \textbf{Требование} & \textbf{Способ проверки в работе} \\
        \hline
        Сохранность принятого события & Создание события фиксируется в Facade PostgreSQL вместе с outbox-записью; дальнейшая публикация выполняется из сохранённого состояния. \\
        \hline
        Идемпотентность входной команды & Повторный вызов CreateNotificationEvent с тем же idempotency key не должен создавать новое событие. \\
        \hline
        Контролируемая повторная обработка Kafka-сообщения & Повторное получение сообщения sender-сервисом должно приводить к чтению уже обработанной записи или к подавлению дубля через inbox и уникальные ограничения. \\
        \hline
        Наблюдаемость маршрута & Маршрут подтверждается метриками gRPC и Kafka, а также записью итогового статуса в history-writer. \\
        \hline
        Мультиканальность & Один и тот же базовый процесс проверяется на двух канальных контурах: notification-mail-sender и notification-sms-sender. \\
        \hline
    \end{tabular}
\end{table}

\subsection{Критерии сравнения архитектурных вариантов}

Архитектурные варианты сравнивались не по общему удобству разработки, а по критериям, значимым для платформы уведомлений:
\begin{enumerate}
    \item сохранность принятого события при частичных отказах;
    \item возможность вынести фактическую доставку из критического пути клиентского запроса;
    \item поддержка повторной обработки без неконтролируемых дублей;
    \item независимое масштабирование приёма команд, планирования и канальных обработчиков;
    \item возможность подключения нового канала без переработки базовой модели;
    \item сложность эксплуатации и диагностики;
    \item сложность начальной реализации.
\end{enumerate}

\subsection{Сравнение архитектурных вариантов}

\begin{table}[H]
    \centering
    \footnotesize
    \renewcommand{\arraystretch}{1.15}
    \caption{Сравнение архитектурных вариантов платформы уведомлений}
    \label{tab:architecture-variants}
    \begin{tabular}{|L{0.24\textwidth}|L{0.29\textwidth}|L{0.31\textwidth}|}
        \hline
        \textbf{Вариант} & \textbf{Преимущества} & \textbf{Ограничения} \\
        \hline
        Монолит с канальными адаптерами & Простая начальная реализация, единый процесс и единая точка отладки & Не выносит доставку из критического пути, не даёт независимого масштабирования каналов и объединяет все отказы в одном процессе \\
        \hline
        Модульный монолит & Улучшает внутреннюю структуру кода и снижает хаотичное смешение ответственности & Сохраняет общую runtime-границу, общий релизный цикл и не отделяет канальную обработку от приёма команд \\
        \hline
        Синхронные микросервисы & Разделяют кодовую базу по сервисным ролям и позволяют изолировать часть бизнес-логики & Сохраняют длинный RPC-маршрут, хуже переносят временные отказы downstream и усложняют безопасную повторную обработку \\
        \hline
        Комбинированный микросервисный подход: gRPC для команд, Kafka для фоновой обработки & Отделяет командный контур от доставки, допускает outbox/inbox, независимое масштабирование sender-сервисов и трассируемую историю статусов & Имеет более высокую начальную сложность и требует дисциплины в телеметрии, идемпотентности и обработке повторов \\
        \hline
    \end{tabular}
\end{table}

Монолит с канальными адаптерами не удовлетворяет сразу нескольким критериям из предыдущего раздела. Он не выносит фактическую доставку из критического пути клиентского запроса, не даёт независимого масштабирования каналов и делает сохранность события зависимой от состояния одного процесса. Такой вариант пригоден для малого числа сценариев, но плохо соответствует задаче управляемой повторной обработки и трассируемой истории.

Модульный монолит улучшает структуру кода, однако по ключевым для работы критериям остаётся промежуточным решением. Он удобнее с точки зрения сопровождения, но приём команд, планирование и канальная отправка по-прежнему делят общую runtime-границу и общий релиз. Поэтому критерии независимого масштабирования и изоляции фонового контура в нём выполняются лишь частично.

Синхронные микросервисы удовлетворяют критерию сервисного разделения лучше, чем оба предыдущих варианта, но проигрывают по устойчивости маршрута. Если sender-сервис, profile-consent или внешний провайдер отвечает медленно либо возвращает ошибку, это немедленно отражается на всей цепочке RPC-вызовов. В таком варианте сложнее вынести доставку из пользовательского запроса, а значит хуже выполняются критерии устойчивости к частичным отказам и контролируемой повторной обработки.

В работе выбран комбинированный микросервисный подход: синхронный gRPC-контур используется для команд и запросов состояния, а Kafka --- для фоновой обработки и межсервисной передачи запусков доставки. Этот вариант лучше остальных удовлетворяет прикладным критериям работы. Он позволяет сохранить событие локальной транзакцией, вынести фактическую доставку из клиентского запроса, независимо масштабировать scheduler и sender-сервисы, а также удерживать повторную обработку в контролируемых границах за счёт outbox, inbox, idempotency и истории статусов. По этой причине дальнейшее проектирование и реализация опираются именно на него \cite{newman_microservices,richardson_patterns}.

\subsection{Сравнение транспортов межсервисного взаимодействия}

Сопоставление транспортов приведено в таблице~\ref{tab:transport-compare}.

\begin{table}[H]
    \centering
    \footnotesize
    \renewcommand{\arraystretch}{1.15}
    \caption{Сравнение транспортов межсервисного взаимодействия}
    \label{tab:transport-compare}
    \begin{tabular}{|L{0.16\textwidth}|L{0.25\textwidth}|L{0.24\textwidth}|L{0.25\textwidth}|}
        \hline
        \textbf{Транспорт} & \textbf{Сильные стороны} & \textbf{Ограничения} & \textbf{Роль в работе} \\
        \hline
        gRPC & Строгие protobuf-контракты, низкие накладные расходы, удобная эволюция внутренних API & Не предназначен для массового fan-out и длительных фоновых цепочек & Основной командный протокол платформы \\
        \hline
        REST & Простая интеграция по HTTP и широкая совместимость & Более слабая типизация и менее строгая контрактность внутренних вызовов & Рассматривался как внешний альтернативный вариант \\
        \hline
        Kafka & Независимое потребление, буферизация нагрузки, естественная модель событийного обмена & Требует прикладной идемпотентности и контроля повторной обработки & Основной транспорт фоновой обработки \\
        \hline
        RabbitMQ & Удобен для очередей задач и маршрутизации сообщений & Менее естественен для долгоживущего событийного потока и повторного чтения & Рассматривался как альтернатива, не выбран \\
        \hline
    \end{tabular}
\end{table}

gRPC выбран для командного контура потому, что он лучше остальных транспортов удовлетворяет критериям контрактности и предсказуемого синхронного взаимодействия. Для операций CreateNotificationEvent, TriggerDispatch и чтения истории нужен немедленный ответ, строгая схема данных и совместимая эволюция внутренних API. Эти требования gRPC закрывает лучше, чем REST, поскольку protobuf-контракт строже, а накладные расходы внутреннего вызова ниже.

Kafka выбран для фонового контура по другой группе критериев. Здесь требуется вынести доставку из клиентского запроса, переживать всплески нагрузки, разделить консьюмеров по каналам и допустить повторное получение сообщения брокером без потери управляемости. По этим критериям событийный поток подходит лучше, чем синхронный RPC, поскольку канал доставки перестаёт блокировать командный путь и может масштабироваться отдельно.

REST в рамках платформы уступает gRPC как внутренний командный транспорт, поскольку при том же синхронном характере взаимодействия даёт более слабую контрактность и менее строгую схему типов. Он допустим как внешний интеграционный стиль, но не улучшает ни сохранность события, ни управляемость повторной обработки, ни разделение командного и фонового контуров.

RabbitMQ удобен для очередей задач, но в данной работе требуется не только доставка задания исполнителю, а воспроизводимый событийный маршрут с несколькими типами сообщений, разными консьюмерами и последующей статусной обработкой. По критериям повторного чтения событий, разделения потоков и наблюдаемости маршрута Kafka лучше соответствует выбранной архитектуре \cite{grpc_docs,kafka_docs,kafka_in_action}.

\subsection{Варианты работы с профилем получателя}

\begin{table}[H]
    \centering
    \footnotesize
    \renewcommand{\arraystretch}{1.15}
    \caption{Сравнение вариантов работы с профилем получателя}
    \label{tab:profile-options}
    \begin{tabular}{|L{0.24\textwidth}|L{0.28\textwidth}|L{0.28\textwidth}|}
        \hline
        \textbf{Вариант} & \textbf{Преимущества} & \textbf{Недостатки} \\
        \hline
        Прямой доступ sender-сервисов к общему хранилищу профилей & Минимум промежуточных сетевых вызовов & Нарушает владение данными, дублирует логику правил и затрудняет эволюцию модели \\
        \hline
        Локальная копия профиля в каждом обработчике & Быстрая локальная проверка & Сложная синхронизация и риск устаревания данных \\
        \hline
        Выделенный сервис profile-consent & Единая трактовка правил, единый контракт и централизованное развитие модели & Дополнительный сетевой hop и зависимость от доступности сервиса \\
        \hline
    \end{tabular}
\end{table}

Прямой доступ sender-сервисов к общему хранилищу даёт минимальную сетевую задержку, но не удовлетворяет критерию владения данными: правила допустимости доставки начинают жить не в одном контракте, а во всех канальных обработчиках сразу. В результате изменение профиля или бизнес-ограничений требует менять не один сервис, а сразу несколько контуров.

Локальная копия профиля формально ускоряет чтение, но не удовлетворяет критерию актуальности данных. Цена такого решения --- синхронизация копий, риск устаревших данных и дополнительная логика распространения обновлений. Для профиля получателя, который участвует в решении об отправке, такой риск нежелателен.

Выделенный сервис profile-consent выбран потому, что именно он одновременно удовлетворяет критериям единообразия правил, сервисного владения данными и управляемой эволюции модели. Дополнительный сетевой hop в этом случае является осознанной платой за консистентность трактовки ограничений и за отсутствие прямого чтения Redis из нескольких контуров \cite{newman_microservices}.

\subsection{Сценарии использования}

\textbf{Сценарий 1. Создание уведомительного события.}
Клиентская система передаёт параметры события и аудиторию. Платформа валидирует запрос, проверяет ключ идемпотентности и подтверждает приём команды.

\textbf{Сценарий 2. Немедленный запуск доставки.}
Для выбранного события формируется запуск доставки. Платформа публикует его в рабочий Kafka-топик и завершает синхронный запрос.

\textbf{Сценарий 3. Отложенный запуск доставки.}
Если в запросе указано будущее время отправки, запуск сначала сохраняется в контуре планировщика, а после наступления заданного времени публикуется в рабочий топик.

\textbf{Сценарий 4. Канальная обработка.}
Mail- или SMS-sender принимает запуск, запрашивает решение у profile-consent, выполняет отправку или фиксирует обоснованный пропуск.

\textbf{Сценарий 5. История доставки.}
Статусное событие поступает в history-writer, где сохраняется факт изменения статуса и становится доступной агрегированная сводка по получателю или событию.
